\documentclass{if-beamer}

% --------------------------------------------------- %
%                  Presentation info	              %
% --------------------------------------------------- %
\title[Lecture 5]{Lecture 5}
\subtitle{Variables, Operators, Libraries, Input/output}
\author{Instructor: Ashley Gannon}
\date{ISC3313 Fall 2021}
\logo{
\includegraphics[scale=0.08]{figures/FSULogo.png}
}
\subject{Presentation subject} % metadata

\graphicspath{{figures/}}
% --------------------------------------------------- %
%                    Title + Schedule                 %
% --------------------------------------------------- %
\begin{document}

\begin{frame}
  \titlepage
\end{frame}
% --------------------------------------------------- %
%                      Presentation                   %
% --------------------------------------------------- %
\section{Variables}

\begin{frame}
\frametitle{What is a variable?}
\begin{block}{Short Answer} It is an object that represents a certain class of data \end{block}
For example if you saw the expression $x + 5$ in a math class you would immediately know that $x$ was a variable that represented a number. \\
\medskip
In order to perform any kind of algorithm you must have some sort of data and that data is stored in what we call a variable. \\
\medskip
The class or group of data that the variable might contain is called its type. \\
\medskip
It might be useful to portray data stored in our computers in the following fashion:
\begin{table}
	\begin{tabular}{cl l l}
		\textbf{Variable Name} & \textbf{Type} & \textbf{Value}\\
		x & int & 2 \\
		y & int & 3 \\	
	\end{tabular}
\end{table}
\end{frame}

\begin{frame}
\frametitle{Types}
A few built in types that are available for you to use:
\begin{itemize}
\item \textbf{int} (represents the math class integer with bounds [-2147483648, 2147483647])*
\item \textbf{float} (represents the math class Real with precision $~10^{-7}$)*
\item \textbf{char} (represents a single character of text ie. 'a')
\item \textbf{string} (We've already seen this one in our Hello World example. It represents a list of characters.)
\end{itemize}
\medskip
*Values may vary from one computer to another.
\end{frame}

\begin{frame}
\frametitle{Variable Declaration and Assignment}
To declare a variable we state the type then the name. \\
For example if we typed \texttt{int x;} we would create a variable named x of type int 
\begin{table}
\begin{tabular}{cl l l}
\textbf{Variable Name} & \textbf{Type} & \textbf{Value}\\
x & int & ? \\
\end{tabular}
\end{table}
\end{frame}

\begin{frame}
\frametitle{Variable Declaration and Assignment}
To declare a variable we state the type then the name. \\
For example if we typed \texttt{int x;} we would create a variable named x of type int 
\begin{table}
	\begin{tabular}{cl l l}
		\textbf{Variable Name} & \textbf{Type} & \textbf{Value}\\
		x & int & ? \\
	\end{tabular}
\end{table}

To assign a value to a variable, we use the \texttt{=} operator. For example if we now typed \texttt{x = 5;}, the the variable x would contain the value 5.
\begin{table}
	\begin{tabular}{cl l l}
		\textbf{Variable Name} & \textbf{Type} & \textbf{Value}\\
		x & int & 5 \\
	\end{tabular}
\end{table}
\end{frame}

\begin{frame}
\frametitle{Variable Declaration and Assignment}
To declare a variable we state the type then the name. \\
For example if we typed \texttt{int x;} we would create a variable named x of type int 
\begin{table}
	\begin{tabular}{cl l l}
		\textbf{Variable Name} & \textbf{Type} & \textbf{Value}\\
		x & int & ? \\
	\end{tabular}
\end{table}

To assign a value to a variable, we use the \texttt{=} operator. For example if we now typed \texttt{x = 5;}, the the variable x would contain the value 5.
\begin{table}
	\begin{tabular}{cl l l}
		\textbf{Variable Name} & \textbf{Type} & \textbf{Value}\\
		x & int & 5 \\
	\end{tabular}
\end{table}

It is also possible to declare a variable and initialize it in the same line.
For example: \texttt{int x = 5;}, would give us the same result.
\begin{table}
	\begin{tabular}{cl l l}
		\textbf{Variable Name} & \textbf{Type} & \textbf{Value}\\
		x & int & 5 \\
	\end{tabular}
\end{table}
\end{frame}


\begin{frame}
\frametitle{Exercises}
Try Compiling and running after every step, so that you can spot your mistakes a little easier. 
\begin{itemize}
\item Create a file called Types.cpp
\item Declare and initialize an integer variable named x to the value 5;
\item Declare and initialize a float variable named y to the value 1.5;
\item Delcare an int variable named z;
\item Initialize z to the value $x / y$. 
\item Print the value of z with the command \texttt{printf(z);} Why doesn't it work?
\item Now try printing the value of z with the command \texttt{printf("\%f", z);}
\item What's the value of z? Why do you think that is?
\item What if we change the f to a d? \texttt{printf("\%d", z);}
\end{itemize}
\end{frame}


\begin{frame}
\frametitle{Miscellaneous Notes on Variables}
\begin{minipage}{0.7\textwidth}
\begin{itemize}
\item Variable names are case sensitive. Thus \texttt{patrick} is registered as a different variable than \texttt{Patrick}.
\item When assigning char's you must put a single quote around the character. \\ \texttt{char letterA =  'a';}
\item When assigning a string you must put a double quote around the group of characters, ie. "Hello World!".
\end{itemize}
\end{minipage}
\begin{minipage}{0.3\textwidth}
\includegraphics[width = \textwidth]{figures/patrick}	
\end{minipage}

\end{frame}

\section{}

\begin{frame}
\frametitle{Booleans}
\begin{block}{Boolean}
	A variable that covers the state space of true or false. Which in terms of ones and zeros is either 1 or 0. A small caveat of c++ is that it often treats everything that is not a zero as true. This can be useful or it can be annoying. 
\end{block}
\medskip
A declaration and initialization might look like: \\
\vspace{1pt}
\texttt{bool performCalculation = false;}
\end{frame}

\begin{frame}
\frametitle{Operators}
An operator is a simple algorithm that takes some input (usually two values) and returns another value as output. The mathematical operators are good examples and behave as expected. 
\begin{block}{Comparison operators}
	\vspace{2pt}
	\begin{columns}[c] % The "c" option specifies centered vertical alignment while the "t" option is used for top vertical alignment
		\column{.18\textwidth} % Left column and width
		$<$ less than\\ $>$ greater than 
		\column{.35\textwidth}
		$<=$ less than or equal to\\ $>=$ greater than or equal to 
		\column{.3\textwidth}
		== equal to\\ != not equal to
	\end{columns}
\end{block}
\begin{block}{logical operators}
	\&\& Represents 'AND'\\
	$||$ Represents 'OR' \\
	! Represents 'NOT'
\end{block}
\end{frame}

\begin{frame}
\frametitle{Conditions}
\begin{block}{Condition}
A statement that evaluates to a boolean value. 
\end{block}
\begin{block}{Conditional statements: True or false?}
	\begin{enumerate}
		\item 1 $<$ 2
		\item 2 $>=$ 3
		\item 1 $!=$ 2
		\item 1 $==$ 2
		\item 1 $=$ 2
	\end{enumerate}
\end{block}
\end{frame}

\begin{frame}
\frametitle{Conditions}
\begin{block}{Condition}
	A statement that evaluates to a boolean value. 
\end{block}
\begin{block}{Conditional statements: True or false?}
	\begin{enumerate}
		\item 1 $<$ 2  \hfill  True \\
		\item 2 $>=$ 3  \hfill False \\
		\item 1 $!=$ 2  \hfill True \\
		\item 1 $==$ 2  \hfill False \\
		\item 1 $=$ 2  \hfill Invalid notation \\
	\end{enumerate}
\end{block}
\end{frame}


\begin{frame}
\frametitle{Tokens}
\begin{figure}
	\center
	\includegraphics[width=0.9\textwidth]{tokens.png}
\end{figure}
\end{frame}

\begin{frame}
\frametitle{Keywords}
Some keywords are covered by the datatypes. Some other important
ones to know for now are:
\begin{itemize}
\item \texttt{if}: executes a chunk of code if a given condition is
true (can execute other code with \texttt{else} if condition
is not true)
\item \texttt{for}: allows the program to execute a number of
statements in sequence
\item \texttt{while}: allows the program to execute statements
until a condition is satisfied
\end{itemize}
Much more on these later in the course!
\end{frame}

\begin{frame}
\frametitle{Identifiers}
Identifiers are things that come from outside
the language itself, i.e. from user-written code. For example, the
\texttt{cout} and \texttt{endl} identifiers: \\[0.5cm]
{ \footnotesize 
cout << "Hello world!" << endl;
} \\[0.5cm]
These actually come from a library.
\end{frame}

\begin{frame}
\frametitle{Libraries}
Libraries can be input to a program by means of a \textit{header
file}. This file has a special extension, \texttt{.h}. For example,
the header file \texttt{iostream.h} gives us the \texttt{cout} and
\texttt{endl} identifiers. Some other popular ones are
\begin{itemize}
\item \texttt{stdio.h}
\item \texttt{iostream.h}
\item \texttt{stdlib.h}
\item \texttt{math.h}
\item \texttt{string.h}
\end{itemize}
You can think of these as augmenting the \texttt{C++} language. You
can write your own header files! Put them into your program with
\texttt{\#include <filename.h>}
\end{frame}

\begin{frame}
\frametitle{Input and Output}
Instead of printf();
\medskip

\begin{block}{cout $<<$ variable;}
	Used to print to the terminal. It can be used to print many different data types. 
\end{block}

\begin{block}{cin $>>$ variable;}
	Takes input from the screen and puts it into the following variable. 
\end{block}

These commands are available in the iostream package, so be sure to use \\
\texttt{ \#include $<$iostream$>$} \\
at the top of your program. 

\end{frame}

\begin{frame}
\frametitle{Class Exercise - Loops!}
	Download  \texttt{HelloWorld5\_for.cpp} and \texttt{HelloWorld5\_while.cpp} from Canvas. Move them to your ClassActivity folder. 

\end{frame}



\end{document}
